\section{2007--2019: A Wider Law, and the End of the Commission}

\subsection{\work{Summorum Pontificum} and what it changed for a body that already had the books}

On 7 July 2007 Benedict XVI issued \work{Summorum Pontificum}. Its
central provision, art.~1, states that ``the Roman Missal promulgated by
Pope Paul VI is the ordinary expression of the \latin{lex orandi} \ldots of
the Catholic Church of the Latin rite,'' that ``the Roman Missal
promulgated by Saint Pius V and revised by Blessed John XXIII is
nonetheless to be considered an extraordinary expression of the same
\latin{lex orandi} of the Church and duly honoured for its venerable and
ancient usage,'' and that these two expressions ``will in no way lead to a
division in the Church's \latin{lex credendi} \ldots; for they are two
usages of the one Roman rite.'' Art.~2 permits any Catholic priest of the
Latin rite to use either Missal in Masses celebrated without a
congregation, on any day except the Easter Triduum, ``no permission from
the Apostolic See or from his own Ordinary'' being needed.

Art.~3 is the article addressed to bodies like the Fraternity: ``If
communities of Institutes of Consecrated Life and Societies of Apostolic
Life, whether of pontifical or diocesan right, wish to celebrate the
conventual or community Mass in their own oratories according to the 1962
edition of the Roman Missal, they are permitted to do so. If an individual
community or an entire Institute or Society wishes to have such
celebrations frequently, habitually or permanently, the matter is to be
decided by the Major Superiors according to the norm of law and their
particular laws and statutes.''\endnote{Benedict
XVI, motu proprio \work{Summorum Pontificum}, 7 July 2007, arts.~1--3,
quoted from the Holy See's own English version
(\url{https://www.vatican.va/content/benedict-xvi/en/motu_proprio/documents/hf_ben-xvi_motu-proprio_20070707_summorum-pontificum.html}),
read 2026-07-25; the exact bytes read on that date match the artifact
registered in this repository's source library. Latin in \work{Acta
Apostolicae Sedis} 99 (2007) 777--781.}

For the Fraternity of Saint Peter, \work{Summorum Pontificum} changed
little juridically and a great deal practically. Juridically, the
Fraternity already held the four books by a grant of 1988 that
\work{Summorum Pontificum} neither enlarged nor restricted; what art.~3
added was a general norm assigning to major superiors the decision about
frequent, habitual or permanent community use---which is close to the
power the Holy See had refused the Fraternity's own superiors in 2000,
now conceded by universal law to the superiors of every institute and
society. Practically, the motu proprio removed the argument that had done
the Fraternity most damage. Once the
Holy See taught that the older and newer books are ``two usages of the one
Roman rite,'' the supposition that habitual use of the older constituted a
half-detached ecclesial position lost its footing---and with it, so did
the underlying premise of the 1999 responses, which had reasoned from the
older Missal being an indult standing against a general law. The
Fraternity itself later dated its self-description by that act, and its
own liturgical directory of 2021 speaks of the ``Extraordinary Form of the
Roman Rite'' in \work{Summorum Pontificum}'s
vocabulary.\endnote{Pietras, art.\ cit., 178, citing the Fraternity's
General Directory for the Sacred Liturgy in the \work{Vademecum}
(Fribourg, 2021), nn.~2--5, 19, 24. That directory was not reached for
this account; the point is reported as the article reports it. The
inference in the sentence before is this account's own.}

The instruction \work{Universae Ecclesiae} of 30 April 2011, issued by
the Commission \work{Ecclesia Dei}, applied the motu proprio; the
Fraternity is not named in it, and this account has not analysed
it.\endnote{Congregation for the Doctrine of the Faith---Pontifical
Commission \work{Ecclesia Dei}, instruction \work{Universae Ecclesiae}, 30
April 2011, \work{Acta Apostolicae Sedis} 103 (2011) 413--420; read on
2026-07-25 in the English text on the Holy See's site
(\url{https://www.vatican.va/roman_curia/pontifical_commissions/ecclsdei/documents/rc_com_ecclsdei_doc_20110430_istr-universae-ecclesiae_en.html}).
Cited for its existence and locus; not analysed. The 2011 \work{Acta}
volume is not available in the Holy See's whole-volume PDF archive, and the
locus is taken from the citation in the canonical literature.}

\subsection{The visitation of 2014}

The Fraternity records, in a communiqué of 2024, that ``the last ordinary
apostolic visit of the Fraternity was undertaken in 2014 by the Ecclesia
Dei Commission.'' No decree of appointment, mandate, report or communiqué
of that visitation was located for this account, and nothing is asserted
about it beyond the Fraternity's own statement that it
occurred.\endnote{``Apostolic visitation of the Fraternity'' (fssp.org,
communiqué of 26 September 2024, published 27 September), read 2026-07-25.
Bounded negative: no act of the Pontifical Commission concerning a 2014
visitation was located in \work{Acta Apostolicae Sedis}, on vatican.va, or
in the Holy See Press Office bulletin. The word ``ordinary'' is the
Fraternity's.}

\subsection{17 January 2019: the Commission decommissioned}

The single most important act of this period for the Fraternity's
constitutional position is not liturgical at all. On 17 January 2019
Francis issued an apostolic letter given motu proprio ``on the Pontifical
Commission \work{Ecclesia Dei}.'' Its narrative part rehearses the
Commission's thirty years; its dispositive part is three sentences:

\begin{enumerate}
\item ``The Pontifical Commission \work{Ecclesia Dei}, instituted on 2
July 1988 with the Motu Proprio \work{Ecclesia Dei Adflicta}, is
decommmissioned.''
\item ``The tasks of the Commission in question are assigned entirely to
the Congregation for the Doctrine of the Faith, within which will be
established a special Section committed to continue the work of
supervision, promotion and protection conducted thus far by the
decommissioned Pontifical Commission \work{Ecclesia Dei}.''
\item ``The financial report of the Pontifical Commission returns to the
ordinary accounting of the aforementioned
Congregation.''\endnote{Francis, apostolic letter given motu proprio ``On
the Pontifical Commission Ecclesia Dei,'' 17 January 2019, quoted from the
Holy See's English text
(\url{https://www.vatican.va/content/francesco/en/motu_proprio/documents/papa-francesco-motu-proprio-20190117_ecclesia-dei.html}),
read 2026-07-25. The spelling ``decommmissioned'' is the Holy See's own in
that version. The motu proprio orders its promulgation by publication in
\work{L'Osservatore Romano} of 19 January 2019, immediate entry into
force, and subsequent insertion in \work{Acta Apostolicae Sedis}.}
\end{enumerate}

The reasoning given matters as much as the disposition. The letter states
that the conditions which led John Paul II to institute the Commission
have changed; that ``the Institutes and Religious Communities which
customarily celebrate in the extraordinary form have today found proper
stability of number and of life''; and that the questions the Commission
addressed ``are of a predominantly doctrinal order.'' It records that the
Feria IV of the Congregation for the Doctrine of the Faith of 15 November
2017 asked that dialogue with the Society of Saint Pius X be conducted
directly by that Congregation, and that the pope approved the request on
24 November 2017.

\subsection{The citation that settles the constitutional question}

The letter's second footnote is, for this history, the decisive
sentence in the file. Explaining that the Commission ``was able to
exercise its own authority and competence in the name of the Holy See over
these societies and associations, until otherwise provided,'' it cites:
``\latin{Rescriptum ex Audientia Sanctissimi}, 18 Oct.\ 1988, AAS, LXXXII
(1990), 5 (3 Maii 1990), 533--534, 6.''\endnote{Same motu proprio, note
[2]. The citation is exact: volume 82, fascicle 5 of 3 May 1990, pages
533--534, number 6---that is, the sixth of the faculties granted in
\latin{Quia peculiare munus}, the one qualified \latin{donec aliter
provideatur}. Note [1] cites \work{Ecclesia Dei} 6~a and note [4] cites
\work{Ecclesiae unitatem} 5.}

That is the Holy See, in 2019, identifying by volume, fascicle, page and
number the exact clause under which it was acting. The clause it invokes
is faculty 6: the exercise of the authority of the Holy See over these
societies, \latin{donec aliter provideatur}. Not faculty 3~a, the erection.
Not faculty 3~b, the seminary. Not faculty 1, the Missal. The 2019 act
disposes of the supervisor, and it says so by citing the only faculty in
the 1988 rescript that had been made provisional.

There is no need for this account to argue that the erection survived the
Commission that made it. The act that ended the Commission argues it, by
the number it cites.

\actrecord{Francis, apostolic letter given motu proprio on the Pontifical
Commission \work{Ecclesia Dei}, 17 January 2019.}{That the Commission was
decommissioned; that its tasks passed entire to the Congregation for the
Doctrine of the Faith through a special section; that the stated grounds
included the stability of the institutes celebrating in the extraordinary
form and the predominantly doctrinal character of the remaining
questions; and that the Holy See identified its title to act as n.~6 of
the rescript of 18 October 1988, \work{Acta Apostolicae Sedis} 82 (1990)
533--534.}{An act about a curial organ. It names no institute, alters no
erection, withdraws no liturgical faculty, and says nothing about the
1962 books.}

\begin{documentinventory}
\inventoryrow{\work{Summorum Pontificum}, 7 July 2007}{Holy See English
text; \work{Acta Apostolicae Sedis} 99 (2007) 777--781}{The two-usages
doctrine and the general regime of the 1962 Missal. Confers nothing on
the Fraternity, which already held the books.}
\inventoryrow{\work{Universae Ecclesiae}, 30 April 2011}{Holy See English
text; \work{Acta Apostolicae Sedis} 103 (2011) 413--420}{Cited for
existence and locus only; not analysed.}
\inventoryrow{Apostolic visitation of 2014}{Fraternity communiqué of
2024 only}{That the Fraternity says one occurred. No Roman act located.}
\inventoryrow{Motu proprio of 17 January 2019}{Holy See English
text}{Decommission of the Commission; transfer of its tasks to the
doctrinal congregation; express citation of faculty 6 of the 1988
rescript.}
\end{documentinventory}
